% Abstract: a single paragraph, under 300 words, no citations.

Physical AI is entering oncology clinical trials through in silico simulation,
autonomous code generation, and eventually robotic procedures, settings in which
an error that ships costs far more than one caught before it ships. This work
reports an end to end pipeline in which a single autonomous agent, Claude Code
Opus 4.7 (1M context) Max, performs five operations in order: it generates
bracketed instructions, generates code from them, executes that code, generates
the charts and figures, and then assembles first a draft and finally this full
paper. The central argument is that none of these generation steps decides
whether a deliverable is trustworthy; the verification, validation, and
uncertainty quantification (VVUQ) process does, and it must be held to a higher
standard than the generation it gates. Holding assurance above generation is
what makes autonomous physical AI oncology trial deliverables faster, less
expensive, and more rigorous than conventional verification. The executed
evidence supports the claim directly. The platform, packaged as the
\texttt{cancer-automated} repository, passed 51 of 51 automated tests across 8
modules and accelerated a representative schedule by a factor of 2.5, from 30
baseline days to 12 automated days. Its VVUQ gate, centered on
\texttt{vvuq/vvuq\_gate.py}, was exercised across 6 cases: it accepted 1 and
blocked 5, escalating 1 of those blocks to a human, and it requires a
verification fraction of exactly 1.0 to ship, while the accepted candidate held
a maximum coefficient of variation of 0.0068 against a 0.10 bound. The same gate
scales to a Stage 2 2030 pancreatic cancer analog: a 60 second 8 arm robotic
Whipple followed by six 28 day adjuvant cycles spanning 168 regimen days, each
gated and supervised by a human.
